\chapter{Generalidades}\label{cap1:generalidades}

\section{Introducción}

La simulación de la propagación de la luz es una herramienta fundamental en la
ciencia y la ingeniería, con aplicaciones que abarcan desde el diseño de sistemas
de telecomunicaciones y la nanofotónica hasta el desarrollo de técnicas de
imagenología médica. Un fenómeno de particular interés es el \textit{speckle}
óptico, un patrón de interferencia granular de alto contraste que surge cuando
luz coherente, como la emitida por un láser, se dispersa sobre una superficie
rugosa o un medio difusor.

Tradicionalmente, la simulación de estos fenómenos ondulatorios se aborda
mediante métodos numéricos robustos como el Método de Elementos Finitos
(\textit{FEM}) o métodos estocásticos como el de Monte Carlo. Sin embargo,
estos métodos presentan un cuello de botella computacional significativo:
requieren una discretización suficientemente fina del dominio espacial para
resolver las oscilaciones del campo de luz. Esto puede aumentar el consumo de
memoria y el tiempo de cómputo, especialmente cuando se necesitan muchas
realizaciones estadísticas
\parencite{jin2014finite, goodman2007speckle}.

En paralelo, el campo de la inteligencia artificial ha experimentado un
crecimiento disruptivo. Una innovación reciente son las Redes Neuronales
Informadas por Física (\textit{PINNs}, por sus siglas en inglés
\textit{Physics-Informed Neural Networks}), una metodología que fusiona el
\textit{Deep Learning} con las leyes físicas fundamentales. Las \textit{PINNs}
\parencite{raissi2019physics} son redes neuronales entrenadas no solo para
ajustarse a datos de frontera, sino para obedecer directamente una Ecuación
Diferencial Parcial (EDP) que gobierna el fenómeno físico de interés.

No obstante, las activaciones estándar de las redes neuronales (tanh, ReLU)
presentan un sesgo espectral \parencite{rahaman2019spectral}: tienden a aprender
frecuencias bajas antes que las altas. Para la ecuación de Helmholtz con número
de onda $k = 2\pi$, este sesgo puede dificultar el ajuste de soluciones
oscilatorias y hacer más exigente el cálculo de derivadas de segundo orden.
Las \textit{Sinusoidal Representation Networks (SIREN)}, propuestas por
\textcite{sitzmann2020siren}, utilizan la activación $\phi(z) = \sin(\omega_0 z)$;
por ello se seleccionan en este trabajo como una arquitectura adecuada para
evaluar ecuaciones de onda, sin asumir que sean universalmente superiores.

Este proyecto se sitúa en la intersección de la óptica computacional y la
inteligencia artificial, aplicando el paradigma PINN-SIREN para superar las
limitaciones computacionales de los métodos clásicos en la simulación del
\textit{speckle} óptico.

% ──────────────────────────────────────────────────────────────────────────────
\section{Planteamiento del problema}

\subsection{Definición del problema}

La simulación precisa de cómo la luz se dispersa y forma patrones de
\textit{speckle} es crucial para sistemas ópticos avanzados en imagenología
biomédica (tomografía de coherencia óptica), metrología (medición de
vibraciones y deformaciones) y criptografía óptica.

El desafío central radica en el costo computacional prohibitivo de los métodos
numéricos actuales:

\begin{enumerate}
    \item \textbf{Métodos Basados en Malla (FEM):} Para capturar con precisión
    los detalles finos de un patrón de \textit{speckle}, la malla debe tener
    una resolución espacial menor que la longitud de onda de la luz
    \parencite{jin2014finite}. En dominios grandes o con alta resolución esto
    puede conducir a sistemas con un número muy elevado de incógnitas, además de
    requerir memoria y tiempo de ejecución considerables. Cada nuevo perfil de
    rugosidad superficial puede exigir repetir la resolución del problema.

    \item \textbf{Métodos Estocásticos (Monte Carlo):} Requieren simular un
    número masivo de fotones individuales para construir un patrón estadísticamente
    preciso, lo que los hace prohibitivos para estudios paramétricos que requieren
    miles de realizaciones.
\end{enumerate}

Aunque la aplicación de las \textit{PINNs} para resolver la ecuación de
Helmholtz ha sido validada en simulaciones de ondas acústicas
\parencite{schoder2024helmholtz} y, de forma más reciente, en contextos
bidimensionales de Helmholtz con arquitectura SIREN
\parencite{panagiotakopoulos2026helmholtz}, en la revisión realizada no se
identificó un trabajo que combine simultáneamente la activación sinusoidal,
una regla explícita de calibración de $\omega_0$ respecto al número de onda y
la generación de \textit{speckle} óptico desde una condición de frontera de
fase aleatoria validada estadísticamente conforme al criterio de Goodman. Esta
combinación es la que aborda el presente trabajo y deberá contrastarse con una
revisión sistemática más amplia si se formula como afirmación de novedad.

Este trabajo propone tres elementos que, en conjunto, no se encontraron en la
literatura revisada. Primero, una regla de calibración
$\omega_0 \approx k/(2\pi)$ que vincula el hiperparámetro de frecuencia de
SIREN al número de onda físico, no identificada en la propuesta original de
SIREN \parencite{sitzmann2020siren} ni en
\textcite{panagiotakopoulos2026helmholtz}, el trabajo más cercano
arquitectónicamente encontrado en la revisión de literatura (misma red
SIREN, mismo campo complejo, misma estrategia de optimización Adam seguida
de L-BFGS). Segundo, la generación de \textit{speckle} óptico mediante la
resolución de la ecuación de Helmholtz desde una condición de
frontera de fase aleatoria, en lugar de ajustar un modelo a patrones de
\textit{speckle} ya medidos. Tercero, la validación estadística del patrón
generado conforme al criterio de Goodman, no identificada en los trabajos de
PINN aplicados a la ecuación de Helmholtz incluidos en esta revisión.

La contribución se plantea como un marco de validación reproducible: la PINN se
compara contra una referencia independiente de propagación por espectro angular,
se evalúa con realizaciones de fase no utilizadas durante el ajuste y se
reportan simultáneamente métricas de campo complejo y estadísticas de
\textit{speckle}. La afirmación de novedad queda limitada a esta combinación y
deberá sostenerse únicamente con los resultados que superen los criterios
cuantitativos definidos en el Capítulo~\ref{cap3:modelo}.

\subsection{Delimitación de la investigación}

\subsubsection*{Alcances}
\begin{itemize}
    \item Se entregará un prototipo de \textit{software} funcional, implementado
    en Python (PyTorch), capaz de construir y entrenar un modelo PINN-SIREN para
    resolver la ecuación de Helmholtz en dominios 2D.

    \item Se validará cuantitativamente la precisión del modelo contra soluciones
    analíticas conocidas mediante el error relativo $L^2$ en NB01 y NB02, y contra
    una referencia numérica independiente en NB03.

    \item Se implementará una extensión del modelo para simular la generación
    de patrones de \textit{speckle} óptico en 2D, mediante una pantalla de fase
    aleatoria espacialmente correlacionada y limitada al régimen propagante.

    \item Se evaluará estadísticamente el patrón generado conforme al criterio
    de Goodman para \textit{speckle} completamente desarrollado
    \parencite{goodman2007speckle}, una vez que NB03 alcance una solución física
    consistente.

    \item Se realizará una evaluación comparativa (\textit{benchmark}) del tiempo
    de inferencia contra una referencia numérica de propagación y, cuando el
    problema se formule con fronteras absorbentes, contra una implementación FEM.
\end{itemize}

\subsubsection*{Limitaciones}
\begin{itemize}
    \item \textbf{Dimensionalidad:} Todas las simulaciones se restringen a
    dominios bidimensionales (2D). El desarrollo en 3D está fuera del alcance
    de este trabajo de maestría.

    \item \textbf{Alcance computacional:} Esta es una investigación puramente
    computacional y de simulación. No se incluyen montajes ópticos experimentales
    ni validación contra datos físicos reales.

    \item \textbf{Simplificación del medio difusor:} El medio se modela como
    una frontera con fase aleatoria discreta ($N_\psi = 256$ puntos). No se
    aborda la simulación de medios volumétricos ni propiedades no lineales.
\end{itemize}

% ──────────────────────────────────────────────────────────────────────────────
\section{Preguntas de investigación e hipótesis}

Las preguntas que guían esta investigación son:

\begin{enumerate}
    \item ¿Puede un modelo PINN con activación sinusoidal (SIREN) aproximar
    con precisión la solución de la ecuación de Helmholtz 2D para la simulación
    de patrones de \textit{speckle} óptico, logrando un error complejo relativo
    $L^2 < 5\%$ respecto a una referencia numérica independiente?

    \item ¿Cuál es la ganancia en eficiencia computacional (tiempo de inferencia
    y uso de memoria) del enfoque PINN-SIREN respecto al FEM para simular el
    mismo fenómeno?

    \item ¿Presenta el patrón de intensidad simulado las propiedades estadísticas
    características del \textit{speckle} completamente desarrollado definidas por
    la teoría de Goodman?
\end{enumerate}

\textbf{Hipótesis:} El desarrollo de un modelo PINN-SIREN para la resolución de
la ecuación de Helmholtz 2D con campo complejo permitirá:
(i) simular la generación de \textit{speckle} óptico con un error relativo
    $L^2 < 5\%$ respecto a la referencia correspondiente al experimento:
    analítica en NB01/NB02 y numérica independiente en NB03;
(ii) obtener un factor de aceleración $S = T_\text{FEM} / T_\text{PINN} > 1$ en
la fase de inferencia; y
(iii) preservar las propiedades estadísticas del fenómeno con un contraste de
\textit{speckle} $|C - 1| < 0.1$ conforme al criterio de \textcite{goodman2007speckle}.

% ──────────────────────────────────────────────────────────────────────────────
\section{Objetivo general}

Construir y validar un modelo computacional basado en Redes Neuronales
Informadas por Física con activación sinusoidal (PINN-SIREN) para la simulación
físicamente precisa de patrones de \textit{speckle} óptico, mediante una
formulación modal de la ecuación de Helmholtz bidimensional con campo eléctrico
complejo, en la que la condición de Cauchy se impone por construcción.

\section{Objetivos específicos}

\begin{enumerate}
    \item Diseñar e implementar una arquitectura SIREN que incorpore el residuo
    de la ecuación de Helmholtz en su función de pérdida, con activación
    $\phi(z) = \sin(\omega_0 z)$ y una estrategia de inicialización que preserve
    la distribución de activaciones en profundidad \parencite{sitzmann2020siren}.

    \item Entrenar y validar la precisión del modelo PINN-SIREN en el escenario
    unidimensional (NB01), comparando contra la solución analítica
    $E(x) = \cos(kx)$ mediante el error relativo $L^2$.

    \item Escalar el modelo al caso bidimensional con campo complejo (NB02),
    empleando Muestreo Latino Hipercúbico (LHS) para los puntos de colocación y
    validando contra la solución de onda plana $E = e^{i(k_x x + k_y y)}$.

    \item Aplicar el modelo validado a la simulación de \textit{speckle} óptico
    (NB03) mediante una pantalla de fase aleatoria correlacionada, condiciones
    laterales periódicas y una condición de propagación saliente, y validar el
    campo contra el espectro angular y el patrón conforme al criterio de Goodman.

    \item Comparar cuantitativamente el rendimiento computacional del modelo
    PINN-SIREN contra una implementación FEM de referencia (NB04), cuantificando
    el factor de aceleración $S = T_\text{FEM} / T_\text{PINN}$.
\end{enumerate}

% ──────────────────────────────────────────────────────────────────────────────
\section{Justificación}

Esta investigación se justifica en tres ejes principales:

\begin{itemize}
    \item \textbf{Justificación computacional:} El principal obstáculo en la
    simulación óptica avanzada es el costo computacional de los métodos de
    malla. El enfoque \textit{mesh-free} propuesto evalúa si es posible reducir
    los tiempos de simulación y facilitar estudios estadísticos de \textit{speckle}
    que resultan costosos con los métodos de referencia.

    \item \textbf{Justificación científica:} La simulación precisa de
    \textit{speckle} impacta directamente en la imagenología biomédica
    (tomografía de coherencia óptica), la metrología de precisión y la
    criptografía óptica. Acelerar la simulación tiene consecuencias directas en
    la capacidad de innovar en estas áreas de alto impacto.

    \item \textbf{Justificación programática:} El proyecto se alinea
    estratégicamente con la Línea de Generación y Aplicación del Conocimiento
    (LGAC) en Ciencia de Datos e Inteligencia Artificial de la Maestría en
    Ciencias de la Computación de la UJAT. Aplica una técnica de vanguardia en
    \textit{Deep Learning} para resolver una EDP fundamental de la física,
    demostrando la sinergia entre la IA y la ciencia computacional avanzada.
\end{itemize}

% ──────────────────────────────────────────────────────────────────────────────
\section{Metodología}

La metodología se basa en el marco de trabajo de las \textit{PINNs} propuesto
por \textcite{raissi2019physics}, extendido con la arquitectura de activación
sinusoidal SIREN \parencite{sitzmann2020siren}. El proyecto es un estudio
computacional cuantitativo y experimental, desarrollado en cuatro fases
alineadas con los objetivos específicos.

\begin{description}
    \item[\textbf{Fase 1 - NB01: Validación analítica 1D.}]
    Se implementa el modelo PINN-SIREN ($5 \times 64$ neuronas) para resolver la
    ecuación de Helmholtz 1D con solución analítica conocida $E(x) = \cos(kx)$.
    La función de pérdida combina el error en las condiciones de frontera
    Dirichlet y el residuo de Helmholtz evaluado en puntos de colocación uniformes.
    Se emplea una estrategia de optimización bifásica: Adam (hasta 15{,}000 épocas)
    seguido de L-BFGS (búsqueda de línea \textit{strong Wolfe}).

    \item[\textbf{Fase 2 - NB02: Escala a 2D con campo complejo.}]
    Se escala la arquitectura a $5 \times 128$ neuronas para acomodar el campo
    complejo $(E_\text{real}, E_\text{imag})$. Se introduce el Muestreo Latino
    Hipercúbico (LHS) con $N_c = 3{,}000$ puntos de colocación y $N_b = 300$
    puntos de frontera por lado \parencite{mckay1979lhs}. Se calibra el peso de
    la función de pérdida $\lambda_\text{fís} = 0.1$ mediante ablación.

    \item[\textbf{Fase 3 - NB03: Simulación de \textit{speckle} óptico.}]
    Se aplica el modelo validado a la generación de \textit{speckle} en el dominio
    de propagación $x/\lambda\in[-10,10]$ y $z/\lambda\in[0,20]$. La pantalla
    de fase tiene correlación espacial finita para evitar una condición dominada
    por modos evanescentes. En la dirección transversal se impone periodicidad y
    en la dirección longitudinal se selecciona la propagación saliente mediante
    el campo y la derivada en la interfaz de entrada. Para mejorar la estabilidad,
    el dominio se entrena por bloques cortos en $z$, transfiriendo el campo y su
    derivada entre interfaces. La PINN se compara contra la referencia de
    espectro angular sin usar valores interiores de esa referencia durante el
    ajuste. La validación estadística sigue el criterio de Goodman
    \parencite{goodman2007speckle}: contraste $C = \sigma_I / \langle I \rangle$,
    distribución de intensidad y, como diagnóstico complementario, el test de
    Kolmogorov-Smirnov.

    \item[\textbf{Fase 4 - NB04: Benchmark FEM (trabajo futuro).}]
    Comparación cuantitativa del tiempo de inferencia PINN-SIREN contra una
    implementación de referencia con FEniCSx, para calcular el factor de
    aceleración $S = T_\text{FEM} / T_\text{PINN}$ y la huella de memoria pico.
    Esta fase está programada como trabajo futuro inmediato.
\end{description}
